% !TEX root = ../AImath.v5.tex

\chapter{ 前言}
\section{从"困难重重"到"显而易见"：智能发展的节奏}
回顾人工智能的发展，我们反复看到一种熟悉的节奏：昨天还被认为“困难重重”的问题，今天却变得“顺理成章”。从早期的专家系统到当下的大语言模型，从手工特征工程到端到端的深度学习，这条演进主线从未中断，只是以新的形式不断重现。许多被视为“难题”的问题，往往并非不可解，而是解决方法在当时尚未成熟。

%--整合？
这种由"困难"走向"显而易见"的转变绝非偶然现象，而是研究范式逐步走向成熟的必然结果。当我们真正抓住了问题的核心本质，并为之匹配到恰当的数学工具与计算方法时，那些原本看似不可逾越的障碍就会像冰雪消融般自然化解。

这就像学习骑自行车的过程。初学者觉得平衡与方向控制极其困难，每一个动作都需要刻意思考。但一旦掌握了"重心转移"和"身体协调"的本质要领，骑车就变成了近乎本能的行为。同样，在AI发展中，许多技术突破都源于这种"本质洞察"——当我们找到了描述和解决问题的正确视角，原本复杂的问题就变得清晰可解。

正如牛顿所言：“如果我看得更远，那是因为我站在巨人的肩膀上。”每一次跨越，都是在前人成果之上完成的再组织与再表达。从这个意义上说，我们今天所认为的“显而易见”，或许正孕育着明天的新突破。
%--整合？

\section{数学思维：智能发展的主线}

要理解人工智能的根脉，我们几乎绕不开数学。数学不仅是描述世界的语言，更是推动智能演化的主线。从古希腊的几何学到现代的概率论，从牛顿的微积分到香农的信息论，每一个数学概念的诞生，都为“理解”与“计算”智能提供了新的可能。

%-----这一段待商榷和修改-------
线性代数让我们得以处理高维数据，微积分使得梯度下降成为可能，概率论为不确定性建模奠定了基础，而信息论则揭示了学习与压缩的内在联系。回望人工智能的发展，每一次重大突破几乎都伴随着数学思想的革新或重生。
每一次理论的成熟，都在为“理解智能”铺路。
%----

深度学习的惊人成功，本质上可以看作是这些经典数学思想在现代计算环境下的重新整合与升华。这个现象为什么值得深入思考？因为它揭示了不同数学分支之间的内在统一性。

具体来说，神经网络的反向传播算法完美体现了微积分中的链式法则——通过层层传递误差信号，实现了复杂函数的高效优化；卷积神经网络巧妙继承了信号处理中的卷积理论，将局部连接和权值共享的思想发挥到极致；而近年来大放异彩的注意力机制，则借鉴了概率分布与加权平均的思维，让模型能够动态地关注输入信息中的关键部分。

这种融合不是简单的拼接，而是深层次的数学统一。就像一位技艺精湛的厨师，将不同食材的风味完美融合，创造出全新的味觉体验。深度学习就是AI领域的"美食烹饪"——它将不同的数学思想和谐地结合在一起，产生了远超各部分简单叠加的强大能力。

\section{智能的层次：从符号到学习，再到创造}
回望人工智能的发展史，我们可以看到一条清晰的演化脉络：从以逻辑推理为代表的“符号主义智能”，到以数据驱动为核心的“学习型智能”，再到当下正逐渐显现的“创造性智能”。这不仅是技术范式的更迭，更是人类理解“智能”的方式在不断深化。

最初阶段的人工智能以符号与规则推理为基础。二十世纪中叶的第一代研究者相信，只要能把人类的知识形式化，机器就能像人类一样思考。早期的专家系统与定理证明程序正是这一理念的体现——它们擅长“演绎”，却难以处理模糊性与经验的不确定。

人工智能最初的发展阶段建立在符号与规则推理的坚实基础之上。在二十世纪中叶，第一代AI研究者们怀抱着一个简单而有力的信念：只要能够把人类的专家知识成功地形式化为符号规则，机器就能够像人类一样进行逻辑思考。这种思想为什么在当时如此吸引人？因为它提供了一条清晰的路径来模拟人类智能。

早期的专家系统与定理证明程序正是这一理念的典型体现。这些系统在特定领域展现出了令人印象深刻的能力——它们能够严格地遵循逻辑规则进行推理，在数学定理证明、医疗诊断等任务中取得了不错的效果。然而，这些系统也暴露出了明显的局限性：它们虽然擅长"演绎推理"，却难以处理现实世界中普遍存在的模糊性、矛盾和经验的不确定性。

就像一个只会背诵规则的学生，面对没有明确规则的现实问题时就会束手无策。第一代AI系统就像是这样的学生——在规则明确的世界里表现出色，但在需要灵活应对复杂现实的场合就显得力不从心。这种局限性也促使研究者们开始探索新的智能范式。

随着数据与计算能力的增长，人工智能进入了“学习”的时代。机器不再依赖人工制定的规则，而是从海量数据中自动发现模式与规律。这一阶段的智能更接近“模仿”：系统通过神经网络与统计模型，复现人类在语言、视觉和决策上的能力。正如孩子学说话，通过反复尝试与纠错，机器逐渐掌握了理解世界的路径。

当学习能力积累到一定程度，智能便开始显露“创造”的潜质。所谓创造，并非凭空发明，而是在已有知识与模式的基础上重新组合、抽象与泛化，从而生成新的结构与洞见。今天的大语言模型与生成式系统，让我们首次直观地感受到这种“新智能”的雏形：它们不仅能模仿，还能在语义与逻辑的空间中创造出人类未曾明确教给它们的内容。

从符号到学习，再到创造，人工智能的演化实际上映射了人类认知方式的演进——从理性逻辑到经验学习，再到综合创新。真正的“创造”，或许正是理解与模仿的高度统一。这一层层的演化，也构成了我们理解智能的路径。

\section{在浪潮中理解智能}
2025年的今天，几乎每隔几天就会诞生一个人工智能新工具。新的模型、框架和平台层出不穷。
%人们惊叹于AI的速度，也常感到跟不上它的脚步。
技术越快，人越焦虑。技术的指数级进步，让我们体验到前所未有的便利，也带来了同样前所未有的不安——焦虑于被取代、被落下，甚至焦虑于自己理解能力的局限。  

面对这股奔腾的技术浪潮，真正值得思考的问题是：在这一切变化背后，人工智能的本质究竟是什么？我们该如何在更新速度如此惊人的时代，建立一种不被浪潮裹挟的理性理解？  

本书的出发点，是希望与读者一道，在纷繁的技术表象之下，寻找理解人工智能的更深层逻辑。我们并不满足于对算法的罗列与介绍，而是想通过这本书的尝试，揭示其背后的数学结构、历史脉络与思想根基。唯有如此，我们才能真正“看懂”智能的成长线索，而不是被无尽的更新所淹没。  

{\bf 学习的关键，不在于记忆全部，而在于洞见规律。}

正如《一如既往》一书所指出的：“
{\bf 刚开始研究一个领域时，往往感觉有无数的知识点需要记忆。其实并非如此。}
如果能够识别统领这个领域的核心法则——通常是3$\sim$12条——
那么当初看似必须记住的数百万件事物，其实只是这些原则的不同组合。”
理解本质规律，就能在纷繁变化中保持清醒与方向。  

因此，我们并不是“讲授者”，而和大家同行的学习者。  
这本书既是对人工智能数学思维的探索，也是一次共同学习的过程。  
如果在阅读中，你能对某个概念产生新的联想，或在某个例子中看到自己的理解路径——那便是我们写作的意义所在。  

%最大的挑战在于：如何在保持严谨的同时，让内容易于理解。{\bf 数学是一种极为精确的语言，但若缺乏直觉引导，也容易成为理解的障碍。}为此，我们提出以下策略：
%\begin{itemize}
%\item \emph{从直觉出发}：以生活化或熟悉的例子引入，再逐步深入数学细节；
%\item \emph{历史视角}：通过历史演进追溯思想的形成；
%\item \emph{实践结合}：理论与应用相辅相成，让抽象变得可感；
%\item \emph{层次递进}：从简单到复杂，从具体到抽象，循序渐进地搭建知识体系。
%\end{itemize}

带着问题阅读，带着好奇探索——这正是学习的起点，也是智能成长的本源。

\newpage